% ch3 D.1 中段 (原书页 139–144 / PDF p154–p159)

\begin{equation*}
\mathcal{H} \;=\; E_0 \sum_j n_j \;+\; \frac{|X|^{2}}{2} \sum_{j\neq k} b_j^{\dagger}(1-n_j)\cdot T_{jk}\cdot(1-n_k)\, b_k \tag{15}
\end{equation*}

再者，(15) 是完全精确的；“动能限制”在这里已被一个非线性相互作用所取代。

从这里出发继续下去有两条途径。第二种，即直接对角化，我们稍后再来演示；第一种是直接计算玻色子算符的运动方程：

\begin{align*}
i\hbar\,\frac{db_j^{\dagger}}{dt} \;=\; \bigl[\mathcal{H},\, b_j^{\dagger}\bigr] \;=\;& E_0\, b_j^{\dagger} \;+\; |X|^{2} \sum_{k\neq j} b_k^{\dagger}\cdot T_{jk}\,(1-n_j)(1-n_k) \\
&-\; \frac{|X|^{2}}{2} \sum_{k\neq j} \Bigl[(b_j^{\dagger})^{2}\cdot T_{jk}\cdot(1-n_k)\, b_k \;+\; b_k^{\dagger}(1-n_k)\cdot T_{kj}\, n_j\Bigr]
\end{align*}

其中用到了 $[b,\, b^{\dagger}]=1$。

若让这个方程作用于没有激子被激发的基态，最后两项便恒等于零，而第二项中的那些修正项也同样为零。于是它是 $b_j$ 的一个线性方程，可以用线性变换

\begin{equation*}
b_j \;=\; \frac{1}{\sqrt{N}} \sum_k e^{i\mathbf{k}\cdot\mathbf{R}_j}\, \beta_k
\end{equation*}

把它从原子激子变换到行波来对角化。所得的方程是

\begin{equation*}
i\hbar\,\frac{d\beta_k}{dt} \;=\; E_0\, \beta_k \;+\; X^{2}\, T\cdot \beta_k
\end{equation*}

其中 $T$ 是偶极相互作用的 Fourier 变换

\begin{align*}
T_k \;&=\; \sum_{j\neq 0}\, e^{i\mathbf{k}\cdot\mathbf{R}_j}\, T_{0j} \\
&=\; \sum_{j\neq 0}\, \frac{e^{i\mathbf{k}\cdot\mathbf{R}_j}}{R_j^{3}}\, \Bigl\{\, 1 \;-\; 3\,\hat{\mathbf{R}}_{0j}\hat{\mathbf{R}}_{0j} \,\Bigr\}
\end{align*}

注意 $\sum_k T_k = 0$，因为其中不含 $j=0$ 的项。对于波长非常长的激子——也只有这类激子才真正值得注意——并且在立方晶体中或沿对称轴的方向上，$T$ 通常可以通过把纵波与横波分开而对角化；也就是说，它的一个主轴沿 $\mathbf{k}$ 方向，另外两个与之垂直。

在 $\mathbf{k}\to 0$ 的极限下，$T$ 变成著名的偶极和（dipolar sum），它在局域场理论中具有基本的重要性。当 $\mathbf{k}\to 0$ 时，

\begin{equation*}
(T_l)_k \;-\; (T_t)_k \;\to\; 4\pi N
\end{equation*}

而对立方对称的格位 $j$，

\begin{equation*}
(T_l)_k \;\to\; \frac{8\pi N}{3} \qquad\qquad (T_t)_k \;\to\; -\,\frac{4\pi N}{3}
\end{equation*}

$4\pi/3$ 正是适当的 Lorentz 局域场因子。$\mathbf{k}\to 0$ 时纵波与横波之间的差别反映了偶极力的长程性质：纵向偶极波会在波的波节与波腹处引起电荷的积累（译注：原文此处作 “modes”，据图 44 及上下文应为 “nodes”，即波节），而横波则不会（图 44）。

%FIG{44}: 图 44　纵向偶极波中偶极子（箭头）与电荷（+、−）的排列：+ 电荷一行处于波节（node），− 电荷一行处于波腹（antinode）（原书图注仅作 “Figure 44”）

于是在这个近似下，激子的能量为

\begin{gather*}
E_l \;=\; E_0 \;+\; \frac{8\pi}{3}\, N X^{2} \;+\; \text{const.}\; X^{2}\, \frac{k^{2}}{a} \;+\; \cdots \\
E_t \;=\; E_0 \;-\; \frac{4\pi}{3}\, N X^{2} \;+\; \text{const.}\; X^{2}\, \frac{k^{2}}{a} \;+\; \cdots
\end{gather*}

关于偶极和，最好的处理大概要数 Cohen 与 Keffer 的一篇文章 (84)。

$T_l$ 与 $T_t$ 的实际数值依赖于结构的具体细节。但二者的差——其实也就是纵向与横向激子频率之差——即使在 Frenkel 模型并不适用的情形下也不依赖这些细节，而是一种纯粹长程的、宏观的效应。各位可以看到，$T_t$ 是由平行的偶极子产生的场之和，其中在极远距离处的表面电荷或体电荷实际上已相互抵消：

%FIG{dipole-t}: 无编号插图（原书不编号）　$T_t$ 的等效图像：左为椭球内交替反向的平行偶极子阵列，椭球表面上 +、− 电荷交替排列；右侧旁注 “which is equivalent to”（它等价于）一个充电的平行板电容器，下方旁注 “or”（或），即一根两端各带 −、+ 电荷、内部箭头同向的细棒

而 $T_l$ 则是相反的情形，等价于

%FIG{dipole-l}: 无编号插图（原书不编号）　$T_l$ 的等效图像：圆柱内彼此平行的偶极子（箭头一律朝上），上表面排满 + 电荷，下表面排满 − 电荷

如图所示。这一差别完全由宏观的表面电荷或体电荷造成，后者可以用宏观介电理论算出。它与 Lyddane、Sachs 与 Teller (85) 等人关于光学声子纵–横频率差的理论完全类似。

为了表明这一切与经典介电理论的直接关系，让我们写下经典电介质的运动方程：

\begin{equation*}
\frac{d^{2}\mathbf{P}}{dt^{2}} \;+\; \omega_0^{2}\,\mathbf{P} \;=\; \beta\, \mathbf{E}_{\mathrm{eff}}
\end{equation*}

这里 $\beta$ 是一个耦合常数，考虑零频率极限便可以用偶极矩阵元把它求出来；在这个极限下，我们从直截了当的微扰论知道，原子系统的极化率为

\begin{equation*}
\alpha \;=\; \frac{2 N X^{2}}{E_0} \;=\; \frac{2 N X^{2}}{\hbar\omega_0} \;=\; \frac{\beta}{\omega_0^{2}}
\end{equation*}

在局域场近似下，原子所感受到的有效场为 $\mathbf{E}_{\mathrm{eff}} = \mathbf{E} + L\mathbf{P}$，其中 $L$ 是“局域场常数”（local-field constant），对立方晶格与横波它取 $4\pi/3$。于是共振频率由

\begin{equation*}
\left(-\omega_r^{2} + \omega_0^{2}\right)\mathbf{P} \;=\; \beta\, L\mathbf{P}
\end{equation*}

决定，并且取 $\omega_r \approx \omega_0$，就得到

\begin{align*}
2\,(\omega_r - \omega_0)\,\omega_0 \;&=\; -\,\beta L \\
\hbar\omega_r \;\simeq\; \hbar\omega_0 \;-\; \frac{\hbar\beta}{2\omega_0}\, L \;&=\; \hbar\omega_0 \;-\; N X^{2} L
\end{align*}

这跟我们用直接的量子力学手段得到的表达式相同。后面我们会看到如何恢复与更完整的经典公式

\begin{equation*}
\frac{\omega^{2} - \omega_0^{2}}{\omega_0^{2}} \;=\; -\,\alpha(\omega{=}0)\cdot L \tag{15'}
\end{equation*}

的相符。

Hopfield 的论文里有一个对整个激子问题的有趣处理，即把一个经典介质作量子化；我不打算在此重述，因为后面讨论自旋波时我们会给出一个类似的处理。

不过，Hopfield 关于激子与辐射相互作用的结果我还是应当讨论一下，因为它们说明了若干有趣的物理要点。到目前为止，我们在计算中只纳入了电磁场中导致 Coulomb 相互作用的纵向部分。横向辐射场与纯纵向的激子并不相互作用，但它当然会通过如下形式的项直接与横向激子耦合：

\begin{equation*}
\sum_{k,\lambda} \left(\mathbf{M}\cdot\mathbf{E}\right)_{k} \left(-\, b_{k\lambda}^{\dagger}\, a_{k\lambda} \;+\; b_{k\lambda}\, a_{k\lambda}^{\dagger}\right)
\end{equation*}

其中 $a_{k\lambda}$ 是光子的产生与湮灭算符，而矩阵元可以证明为

\begin{equation*}
i\left(\frac{2\pi\hbar c}{|k|}\right)^{1/2} \frac{\omega_0}{c}\, X
\end{equation*}

$\lambda$ 是偏振指标。Hopfield 指出，在理想晶体中，我们不应当把这个耦合看作导致光子被激子所吸收，而应当把它看作两个独立的玻色型激发之间的耦合，二者的谱如图 45 所示。激子的速度远远小于

%FIG{45}: 图 45　光子（直线，斜率 = c）与激子（自 $\omega_r$ 出发的缓升曲线）的色散谱，$\omega$ 轴上标出 $\omega_r$；虚线示引入耦合后谱在交叉点附近发生的劈裂（原书图注仅作 “Figure 45”）

光子的速度，因此两支频率在非常小的 $k$ 处相交。自然，当我们引入耦合后，结果便是在交叉点附近两支谱发生劈裂，如图中虚线所示。于是，在没有声子与缺陷的情形下，光子并不被吸收，而是在激子频率处受到强烈色散。对能量处于 $\omega_r$ 附近的辐射，存在一个完全不透明的区域，即截止带（stop band）；在该带的两侧，光子速度被大大减慢。这正是辐射的共振俘获（resonance trapping）现象在固体中的对应物；可以把它想成光子经由反复的吸收与再发射而被减慢乃至停止。当然，在实际晶体中，结果表明通过多声子发射这类过程，激子被吸收的程度远比光子强烈，因而在激子频率附近会发生一个类似吸收的过程。然而必须把这一点理解为从激子-光子混合模式向声子或杂质自由度的直接吸收，而不是光子被激子本身所吸收。

Hopfield 与 Thomas 已经在 CdS 晶体中用实验演示了激子的若干这类性质。通过观察主轴以外方向上的吸收，他们得以造成纵向与横向激子的轻微混合，从而在实验上证实了这两类模式之间的频率差。在另一些安排中，他们借助磁场证明：与光子发生强相互作用的激子确实是以有限的、非零的动量在传播，而不是静止的受激原子；并已在实验上粗略测得激子的速度 (86)。

Frenkel 激子的另一个方面，是高能量偶极项所起的作用；这些项在我们迄今为止的处理中一直被忽略。静电相互作用 (11) 也包含如下形式的偶极矩阵元：

\begin{equation*}
\left(c_j^{p_x}\right)^{\dagger} \left(c_k^{p_x}\right)^{\dagger} c_k^{s}\, c_j^{s} \left(\int a_{p_x}^{*}\, ex\, a_s\right)^{2} \left(T_{jk}\right)_{xx}
\end{equation*}

它们的作用是把原子 $j$ 与原子 $k$ 同时激发。由此在哈密顿量中产生的项，用玻色子 $b$ 写出即为

\begin{gather*}
\frac{X^{2}}{2} \sum_{j\neq k} b_j^{\dagger}(1-n_j)\cdot T_{jk}\cdot b_k^{\dagger}(1-n_k) \\
+\; \frac{X^{*2}}{2} \sum_{j\neq k} (1-n_j)\, b_j\cdot T_{jk}\cdot(1-n_k)\, b_k \tag{16}
\end{gather*}

在大多数情形下我们可以取 $X$ 为实数。（译注：上式第二行的求和限原书印作 $i\neq k$，显系 $j\neq k$ 之误排，此处径改。）

眼下最显而易见的做法，是用微扰论来处理这些项，同时十分有理由地假定相应激发态的平均能量恰好就是 $2E_0$。我们得到固体总能量降低了

\begin{equation*}
\Delta E_2 \;=\; -\, \frac{N|X|^{4}}{2E_0} \sum_j T_{0j}:T_{0j}
\end{equation*}

它正比于 $R_{0j}^{-6}$，显然恰恰应当把它等同于原子间 Van der Waals 吸引中来自这一特定电子跃迁的分量——也就是按独立原子计算出来的那一份。

用另一种方式来做这件事是有趣的。我们看到，在用较简单的哈密顿量 (15) 考虑\emph{第一个}激子时，激子算符运动方程中的非线性项恰好为零。因此在那一级近似下，第一个激子是系统的一个\emph{精确}激发态。现在我们做出典型的“元激发”假设：即使存在物理上数量可观的元激发，元激发之间的相互作用也可以在某级近似下予以忽略；也就是说，无论在 (15) 中还是在我们现在所考虑的 (16) 的各项中，都把 $(1-n_j)$ 换成 1，只保留玻色子振幅的二次项。

必须认识到，一旦把 (16) 计入哈密顿量，就连基态也会含有原子激发态的小振幅，因为存在一些被激发而又退激发的激子对。这个百分比实际上并不

% ——上句在原书 p144 末中断于 “The percentage is not in fact”，p145（PDF p160）起由后续代理续译衔接。
